\answer{ex:20:1}{$t_{\text{v}}=\tau_r\ln\frac{1-L}{1-H}=16.8$~h, $t_{\text{s}}=\tau_d\ln\frac HL=5.8$~h, periodo $22.5$~h; hasta $24$~h lleva al oscilador la oscilación diaria de los umbrales: enganche de fase~\eqref{eq:pll}.}
\answer{ex:20:2}{(a)~$L=0.111$; con $L=0.092$ el sueño habría durado $9.6$~h. (b)~Un $22\%$ ($1/0.82=1.22$), no un $18\%$.}
\answer{ex:20:3}{$H=\frac{p-1}{p-1/q}=0.623$, $L=H/q=0.093$, donde $p=e^{16/\tau_r}$, $q=e^{8/\tau_d}$. Tras despertarla a las $4$~h, la fase se desplaza para siempre (se duerme $7.2$~h antes); con umbrales oscilantes vuelve en dos o tres días.}
\answer{ex:20:4}{(a)~$r_\pm=\frac12\bigl(1\pm\sqrt{1-4k/w}\bigr)=0.947,\ 0.053$; $a_{\text{sup}}=3.51$, $a_{\text{inf}}=0.49$. (b)~Actividad $\approx0.33$~s, silencio $\approx0.59$~s, periodo $\approx0.93$~s, fracción de actividad $0.36$.}
\answer{ex:20:5}{$a_{\text{sup}}/r_+<b<a_{\text{inf}}/r_-$: $3.71<b<9.20$. \nt{ACET} baja $b$ por debajo de $3.71$: la intersección pasa a la rama superior, y el estado activo ($r\approx1$) es estable.}
\answer{ex:20:6}{$4.9$, $6.8$, $8.8$, $5.5$, $8.9$~h con $D=24$, $32$, $40$, $48$, $64$: la duración la fija la fase diaria en que se duerme, no la deuda.}
\answer{ex:20:7}{Tras B, el error en A es en promedio $0.51$ (de $0.19$ a $1.08$ en $20$ conjuntos; con respuesta nula sería $1$). Intercalando, ambos errores son menores que $10^{-4}$.}
\answer{ex:20:8}{Ejercicio abierto. El escalado uniforme conserva los cocientes de los pesos, pero solo conserva las respuestas de la neurona con umbral si el umbral se escala en el mismo factor; las repeticiones intercaladas cambian los cocientes. Que los contactos grandes no cambien contradice un escalado puramente uniforme.}
